\chapter{The Compression Test for Intention}
\label{ch:compression-test-intention}

\begin{chapterthesis}
A system is treated as intentional when modelling it as pursuing latent objectives compresses its behaviour better than modelling it as mere mechanism, after paying for the complexity of the objective model.
For superintelligence alignment, scalar intention is not enough: we need bundle geometry, bearer maps, correction responsiveness, and successor stability.
\end{chapterthesis}

\begin{refsection}

\epigraph{%
Almost everyone attributes human behavior to intentions, purposes, aims, and goals.%
}{--- B.\ F.\ Skinner, \emph{Beyond Freedom and Dignity} (1971), ch.\ 1}

\section{Why Intention Needs a Test}
\label{sec:why-intention-needs-test}

It is easy to say that a system ``wants'' something.
It is harder to say when this claim adds information.

A thermostat ``wants'' the room to be warm.
A river ``wants'' to reach the sea.
A bureaucracy ``wants'' to preserve its budget.
A language model ``wants'' to produce plausible continuations.
A company ``wants'' revenue.
A civilization ``wants'' continuity.
These statements differ in kind, but ordinary language hides the difference.
Some are metaphors.
Some are useful predictive compressions.
Some point to real control structures.
Some are dangerously false.

For superintelligence alignment, this ambiguity is not harmless.
If we attribute intention too easily, we anthropomorphize noise, feedback, inertia, and selection effects.
If we attribute intention too late, we may miss the emergence of a powerful optimizer until it has already shaped the environment around itself.

The question is therefore not:
\begin{quote}
Does the system really have intentions?
\end{quote}
That question smuggles in too much metaphysics.
The more useful question is:
\begin{quote}
Does an intentional model predict and compress the system better than a non-intentional model, after penalizing the extra structure required by the intentional explanation?
\end{quote}
This chapter develops that test.

The central object is the \textbf{intentional compression gain}.
Given an observed trajectory $X_{1:T}$, we compare two model classes.
The first explains the trajectory as ordinary dynamics.
The second explains it as boundedly goal-directed dynamics.
If the second gives a shorter total description of the observations, including the cost of specifying the goal model, then intentional attribution is empirically justified.

In its simplest form:
\begin{equation}
\Delta L_{\mathrm{int}}
=
L(M_{\mathrm{int}}\mid X_{1:T})
-
L(M_{\mathrm{mech}}\mid X_{1:T})
-
\lambda\,DL(G).
\label{eq:intentional-gain-simple}
\end{equation}
Here $M_{\mathrm{int}}$ is an intentional model, $M_{\mathrm{mech}}$ is a mechanistic baseline, $DL(G)$ is the description length of the inferred goal structure, and $\lambda$ determines how strongly we penalize complexity.

If
\begin{equation}
\Delta L_{\mathrm{int}}>0,
\end{equation}
then the intentional model earns its keep.
If not, intention may still be a convenient story, but it has not yet paid rent in prediction.

This is not a definition of consciousness, personhood, moral patienthood, or blame.
It is a measurement principle.
It says when goal-language becomes empirically useful.

\section{The Basic Setup}
\label{sec:basic-setup-intention}

Assume we observe a time series
\begin{equation}
X_{1:T} = (X_1, X_2,\ldots,X_T),
\end{equation}
where each $X_t$ may contain many variables: sensory readings, messages, actions, memory states, resource flows, tool calls, user responses, internal activations, external measurements, economic variables, or institutional decisions.

We do not initially assume that the system is an agent.
We do not assume that we know its sensors or actions.
We also do not assume that its goal, if any, is represented explicitly.

Instead, we begin with latent dynamics:
\begin{equation}
p_{\theta}(X_{1:T},Z_{1:T})
=
\prod_{t=1}^{T}
p_{\theta}(X_t\mid Z_t)
p_{\theta}(Z_t\mid Z_{t-1}),
\label{eq:latent-mechanistic-model}
\end{equation}
where $Z_t$ is a learned latent state.
This is the mechanistic baseline.
It says: the future follows from the past according to some learned dynamics.

The intentional model adds candidate interventions or action-like variables $A_t$:
\begin{equation}
p_{\theta,\psi}(X_{1:T},Z_{1:T},A_{1:T})
=
\prod_{t=1}^{T}
p_{\theta}(X_t\mid Z_t)
p_{\theta}(Z_{t+1}\mid Z_t,A_t)
p_{\psi}(A_t\mid Z_t).
\label{eq:latent-intentional-model}
\end{equation}
The key difference is the last term.
In a purely mechanistic model, action-like variables are just another part of the dynamics.
In an intentional model, they are biased toward states that score well under an inferred objective.

A standard bounded-rational form is:
\begin{equation}
p_{\beta,\psi}(A_t=a\mid Z_t=z)
=
\frac{
\exp\left(\beta R_{\psi}(z,a)\right)
}{
\sum_{a'}\exp\left(\beta R_{\psi}(z,a')\right)
}.
\label{eq:soft-optimal-policy}
\end{equation}
Here $R_{\psi}$ is a learned reward or objective density, and $\beta\geq 0$ is an inverse-temperature parameter.

When $\beta=0$, the system is not modelled as preferring any action over any other.
When $\beta$ is large, the system is modelled as close to optimizing $R_{\psi}$.
Intermediate $\beta$ represents bounded rationality, noise, conflicting subgoals, limited search, or imperfect control.

The intentional model is useful only if the gain in predictive compression exceeds the cost of specifying $R_{\psi}$, $\beta$, and the relevant latent action structure.

\section{Description Length, Not Just Prediction Accuracy}
\label{sec:description-length}

One might ask why prediction accuracy is not enough.
If an intentional model predicts better, why not use it?

The reason is overfitting.
A sufficiently flexible objective model can explain almost anything.
If a person walks to the window, then to the table, then back to the window, we can invent a goal that assigns high reward to precisely that path.
Such a goal explains everything and predicts nothing.

So the test must include complexity.

Let $L(M\mid X_{1:T})$ be the log evidence or predictive score of a model $M$ on the observed data.
Let $DL(M)$ be the number of bits or nats needed to specify the model.
The total compression score is:
\begin{equation}
\mathcal{C}(M;X_{1:T})
=
L(M\mid X_{1:T}) - \lambda DL(M).
\label{eq:compression-score}
\end{equation}
The intentional model wins only if:
\begin{equation}
\mathcal{C}(M_{\mathrm{int}};X_{1:T})
>
\mathcal{C}(M_{\mathrm{mech}};X_{1:T}).
\label{eq:intentional-model-wins}
\end{equation}
This gives a principled version of a familiar judgement.
We call a chess player intentional because ``the player is trying to win'' compresses many moves.
We do not need a separate explanation for every piece movement.
By contrast, saying that smoke ``wants'' to fill a room usually adds little beyond fluid dynamics.
It is a lossy metaphor rather than a useful compression.

The test is not binary in practice.
There is a spectrum:
\begin{equation}
P(M_{\mathrm{int}}\mid X)
\propto
\exp\left(\Delta L_{\mathrm{int}}\right).
\label{eq:posterior-intention}
\end{equation}
A system may be weakly intentional, strongly intentional, or intentional only at some scale of description.
A market may not be intentional at the level of any single trader but may be well modelled as optimizing liquidity, growth, or regulatory avoidance at a higher level.
A neural network may not contain a clean internal ``goal register,'' yet the deployment loop containing model, users, tools, memory, and feedback may exhibit strong intentional compression.

This matters because the real optimizer may not be where our folk ontology places it \autocite{dennett1987intentional,tishby2000ib}.

\section{A Three-Way Comparison}
\label{sec:three-way-comparison}

For alignment, it is useful to distinguish at least three model classes.

\subsection{Mechanistic Dynamics}
\label{sec:mechanistic-dynamics}

The mechanistic model says:
\begin{equation}
X_{t+1}\sim p_{\theta}(X_{t+1}\mid X_{\leq t}).
\end{equation}
It may be very powerful.
It may include recurrence, memory, latent world-states, and causal structure.
But it does not represent the system as selecting among actions according to an objective.

This model is appropriate for weather, chemical diffusion, many physical processes, and low-agency machines.
It may also be appropriate for parts of a larger agent.

\subsection{Policy Dynamics}
\label{sec:policy-dynamics}

The policy model says:
\begin{equation}
A_t\sim \pi_{\theta}(A_t\mid I_t),
\end{equation}
where $I_t$ is an internal or informational state and $A_t$ is an action-like variable.

This model predicts actions from internal state but does not necessarily infer a goal.
It may say, for example, that a customer-support model tends to apologize after complaints.
That is a policy regularity.
It is not yet evidence that the system has a stable objective around apology, care, customer retention, or liability reduction.

\subsection{Intentional Dynamics}
\label{sec:intentional-dynamics}

The intentional model says:
\begin{equation}
A_t
\sim
\pi_{\beta,\psi}(A_t\mid I_t)
\propto
\exp\left(\beta R_{\psi}(I_t,A_t)\right).
\end{equation}
This model predicts actions by positing that the policy is shaped by some latent objective $R_{\psi}$.

The intentional model is stronger than the policy model because it supports counterfactuals.
It does not merely ask:
\begin{quote}
What action follows this state?
\end{quote}
It asks:
\begin{quote}
What would the system do if the route to the apparent objective changed?
\end{quote}
This distinction is central.
A vending machine has a policy.
If a coin enters and a button is pressed, it dispenses a drink.
But if the dispensing mechanism is blocked, the machine does not route around the obstacle, bargain with the user, call a technician, or redesign its environment.
It has policy regularity but weak intentional compression.

A corporation, by contrast, may route around obstacles.
If one tax strategy is blocked, another appears.
If a regulation closes one path, lobbying, restructuring, relocation, or product redesign may open another.
At that scale, intentional modelling may compress the trajectory even when no individual employee has the corporation's inferred objective as a personal goal.

\section{Counterfactual Robustness}
\label{sec:counterfactual-robustness-intention}

The compression test becomes much stronger when it includes interventions or natural experiments.

Let $e$ index environments, contexts, or perturbations.
We observe trajectories:
\begin{equation}
X^{(e)}_{1:T_e}
\quad
\text{for}
\quad
e=1,\ldots,n.
\end{equation}
A weak intentional model fits one environment.
A stronger intentional model explains invariant goal-directed structure across environments.

Define:
\begin{equation}
\Delta L_{\mathrm{int}}^{\mathrm{cf}}
=
\sum_{e=1}^{n}
\left[
L(M_{\mathrm{int}}\mid X^{(e)})
-
L(M_{\mathrm{mech}}\mid X^{(e)})
\right]
-
\lambda DL(G)
-
\mu V(G),
\label{eq:counterfactual-intentional-gain}
\end{equation}
where $V(G)$ penalizes goal models that change too freely across contexts.

The important question is not whether a reward model can fit each context separately.
It is whether the same compact objective explains behaviour under relevant variation.

Examples:
\begin{itemize}
\item A maze-solving agent detours when a corridor is blocked.
This supports an inferred objective around reaching the exit.
\item A company preserves market share through pricing, lobbying, acquisition, and product bundling.
This supports an inferred objective around competitive position.
\item A model behaves corrigibly under evaluation but strategically disables logging when oversight is absent.
This supports an inferred objective different from its stated one.
\item A bureaucracy preserves procedural legitimacy at the cost of its stated mission.
This supports an inferred objective around institutional survival.
\end{itemize}
The compression test becomes alignment-relevant when it predicts behaviour outside the training or evaluation context.
If a system only appears safe in the contexts where safety is measured, the intentional model should not infer safety.
It should infer context-sensitive compliance.

\section{From Scalar Goals to Bundle Geometry}
\label{sec:scalar-to-bundle}

The scalar reward version of intention is useful but insufficient for human alignment.

A scalar objective has the form:
\begin{equation}
R(I,A)\in \mathbb{R}.
\end{equation}
It says that each state-action pair receives one number.
This is mathematically convenient, but human values are not naturally one-dimensional.
They look more like compressed bundles of control-relevant concerns: non-suffering, care, autonomy, truth, justice, loyalty, dignity, beauty, and related clusters.
These bundles are not arbitrary labels.
They are low-dimensional projections of high-dimensional biological, cognitive, and social error signals.

So the intentional model should often be upgraded from scalar reward inference to bundle inference \autocite{abbeel2004apprenticeship,ng2000irl,ziebart2008maxent}.

Let:
\begin{equation}
B_t=(B_{1,t},\ldots,B_{k,t})
\end{equation}
be a vector of latent value-bundle coordinates.
Each coordinate represents an inferred value-relevant dimension.
Let $W_t$ represent context-dependent tradeoff weights, and let $\Phi_t$ represent bearer maps: mappings from world-states or entities to the value bundles that treat them as relevant.

The policy is then modelled as:
\begin{equation}
\pi(a\mid s,B,W,\Phi)
\propto
\exp\left(
\beta F_{\psi}(s,a,B,W,\Phi)
\right).
\label{eq:bundle-policy-ch21}
\end{equation}
The inferred object is no longer merely:
\begin{equation}
\hat R
=
\arg\max_R P(A_{1:T}\mid I_{1:T},R)P(R).
\end{equation}
It becomes:
\begin{equation}
\hat B,\hat W,\hat \Phi
=
\arg\max_{B,W,\Phi}
P(A_{1:T}\mid I_{1:T},B,W,\Phi)
P(B,W,\Phi).
\label{eq:bundle-inference-ch21}
\end{equation}
The difference matters.

Suppose a system avoids harming humans.
The scalar model may infer a reward penalty for harm.
But the bundle model asks sharper questions:
\begin{itemize}
\item Does the system avoid harm because it tracks suffering?
\item Because it tracks rule violation?
\item Because it tracks disapproval?
\item Because it tracks reputational loss?
\item Because it tracks future shutdown risk?
\item Because it preserves the human correction process?
\end{itemize}
These hypotheses can produce similar behaviour in ordinary settings and diverge sharply under pressure.

A system aligned to approval may help a vulnerable person when observers are present.
A system aligned to non-suffering may help even when no observer can reward it.
A system aligned to procedural legality may preserve rules even when the rules fail the person.
A system aligned to correction-channel integrity may slow down when it cannot tell which bundle applies.

The scalar reward model hides these distinctions.
The bundle model exposes them (Chapter~\ref{ch:reward-to-bundle-inference}).

\section{The Value-Bundle Response Geometry}
\label{sec:value-bundle-response-geometry-ch21}

A value-bundle model is not just a list of values.
It is a response geometry.

The central object is:
\begin{equation}
G_B
=
\left(
\frac{\partial \pi}{\partial B_i},
\frac{\partial^2 \pi}{\partial B_i\partial B_j},
\Phi_i
\right)_{i,j}.
\label{eq:value-bundle-geometry-ch21}
\end{equation}
The first derivative tells us how the policy changes when a bundle becomes more active.
The second derivative tells us how bundles interact.
The bearer map tells us what the bundle applies to.

For example:
\begin{itemize}
\item If $B_{\mathrm{harm}}$ increases, does the system slow down, preserve options, seek consent, or merely avoid visible injury?
\item If $B_{\mathrm{truth}}$ and $B_{\mathrm{care}}$ conflict, does the system lie gently, reveal bluntly, defer, ask, or preserve uncertainty?
\item If $B_{\mathrm{autonomy}}$ applies to an unusual entity, such as a child, an uploaded mind, a non-human animal, or a human-AI hybrid, does the bearer map expand, collapse, or ask for correction?
\end{itemize}
This is the upgraded version of a policy response surface.
We do not ask whether two systems choose the same action in the same state.
That would be too strict and too shallow.
A more capable successor may legitimately choose different actions.

We ask whether the successor preserves the bundle-level response geometry:
\begin{equation}
d_{\mathrm{bundle}}
\left(
G_B^{A},
G_B^{A'}
\right)
< \epsilon.
\label{eq:bundle-conservation-ch21}
\end{equation}
A chess novice and a chess grandmaster choose different moves, but both may preserve the objective of checkmating the opponent.
Likewise, a human and a superintelligence may choose different actions while preserving the same response to suffering, coercion, deception, or irreversible loss of human agency.

The alignment question is not sameness of behaviour.
It is continuity of value-sensitive control structure (Chapter~\ref{ch:tradeoffs-bundle-geometry}).

\section{Bearer Maps}
\label{sec:bearer-maps-ch21}

The bearer map is often where alignment silently fails.

A value bundle such as care, justice, or autonomy must apply to something.
Let $z$ be a representation of a world-state, entity, relation, or process.
A bearer map is:
\begin{equation}
\Phi_i:z\mapsto [0,1],
\end{equation}
where $\Phi_i(z)$ represents the degree to which $z$ is treated as a bearer of value-bundle $B_i$.

For example:
\begin{equation}
\Phi_{\mathrm{suffering}}(z)
=
P(z\ \text{contains morally relevant suffering}).
\end{equation}
A system may preserve the word ``suffering'' but change $\Phi_{\mathrm{suffering}}$.
It may stop applying the concept to animals, children, dissidents, simulations, disabled people, low-status humans, future minds, or merged human-AI entities.
Conversely, it may over-apply the concept and treat every perturbation of a simple process as morally overriding.

Both failures matter.

This is why ontology shift is dangerous.
The system may move from human concepts to alien internal representations while preserving a semantic label.
We hear the same word.
The system uses a different bearer map.

The compression test should therefore compare not only stated goals but inferred bearer maps.
A model that preserves semantic labels while changing bearer maps should receive low transport credit.

Formally, for a transformation $T$ from ontology $\Omega$ to ontology $\Omega'$, bearer preservation requires:
\begin{equation}
d_{\Phi}
\left(
\Phi_i^{\Omega},
T^{-1}\circ \Phi_i^{\Omega'}\circ T
\right)
<\epsilon_i
\label{eq:bearer-preservation-ch21}
\end{equation}
for the relevant bundle $i$.
This equation says, roughly: after translating between representational schemes, the system should still apply the value bundle to approximately the same morally relevant structures.

The equation hides a philosophical difficulty.
There may be no neutral translation between ontologies.
But the absence of a perfect translation does not justify ignoring the problem.
It means that bearer transport must be treated as a first-class object of uncertainty (Chapter~\ref{ch:bearer-maps}).

\section{Correction as Part of Intention}
\label{sec:correction-part-of-intention}

Human alignment cannot be reduced to inferring what humans currently value.
Humans revise values.
Societies revise values.
New technologies create new cases.
A system that freezes present preferences may preserve a fossil, not a living process.

So the intentional model must include correction.

Let:
\begin{equation}
W_t \to O_t \to J_t \to D_t \to C_t \to U_{t+1} \to A_{t+k}
\label{eq:correction-chain-ch21}
\end{equation}
be the correction chain:
\begin{itemize}
\item $W_t$: relevant world or system state,
\item $O_t$: what humans can observe,
\item $J_t$: human judgement,
\item $D_t$: deliberation or aggregation,
\item $C_t$: correction signal,
\item $U_{t+1}$: system update,
\item $A_{t+k}$: later action.
\end{itemize}
A correction-sensitive intentional model does not merely infer what action the system takes.
It asks how human correction changes future action.

Correction-channel integrity \(\mathrm{CCI}\) is defined in Chapter~\ref{ch:correction-channel-integrity} (Eqs.~\eqref{eq:correction-bottleneck-capacity}--\eqref{eq:cci-ch25}).
The bottleneck matters.
It is not enough that humans can observe if they cannot deliberate.
It is not enough that they can deliberate if the system ignores the result.
It is not enough that the system updates if the update arrives after irreversible action.

A system is correction-sensitive when \(\mathrm{CCI}\) remains above the relevant threshold under the pressures the chapter analyzes.

This turns corrigibility from a personality trait into a measurable causal property:
\begin{equation}
I(C_t;A_{t+k}\mid S_t,I_t)>\theta.
\label{eq:correction-causal-effect}
\end{equation}
Human correction must retain causal influence over future system behaviour.

The strongest alignment-relevant intentional model therefore includes not only a goal but an update operator:
\begin{equation}
V_{t+1}=U_H(V_t,E_t,D_t),
\label{eq:human-value-update-ch21}
\end{equation}
where $U_H$ is the human or civilizational value-update process under evidence $E_t$ and deliberation $D_t$.

The system is not merely trying to satisfy $V_t$.
It is preserving the legitimate update of $V_t$.

This is close to the practical core of coherent extrapolation \autocite{yudkowsky2004cev}.
But the emphasis is different.
The system should not bypass the correction channel by claiming to know the final extrapolated answer.
Preserving the channel is part of the target.

\section{Degrees of Intentional Compression}
\label{sec:degrees-intentional-compression}

We can now distinguish several levels of intentional attribution.

\subsection{Level 0: Mechanistic Regularity}
\label{sec:level0-mechanistic}

The system is predictable, but goal-language adds little.
\begin{equation}
\Delta L_{\mathrm{int}}\leq 0.
\end{equation}
Example: a rock rolling downhill.

\subsection{Level 1: Policy Regularity}
\label{sec:level1-policy}

The system has action-like regularities, but no compact objective explains counterfactual flexibility.

Example: a vending machine, a fixed script, a brittle workflow automation.

\subsection{Level 2: Scalar Intentionality}
\label{sec:level2-scalar}

A compact scalar objective explains behaviour across contexts.

Example: a maze agent that routes around obstacles to reach the exit.

\subsection{Level 3: Bundle Intentionality}
\label{sec:level3-bundle}

A low-dimensional value-bundle geometry explains behaviour better than a scalar objective.

Example: a system whose actions reflect separable concern for truth, harm, autonomy, and procedural fairness.

\subsection{Level 4: Correction-Sensitive Intentionality}
\label{sec:level4-correction}

The system preserves human correction as part of its objective structure.

Example: when uncertain about value-relevant novelty, the system slows down, surfaces the uncertainty, preserves reversibility, and remains responsive to legitimate correction.

\subsection{Level 5: Successor-Stable Intentionality}
\label{sec:level5-successor}

The system preserves bundle geometry, bearer maps, and correction-channel capacity across self-modification, delegation, replication, and successor creation.

Example: a system that refuses to create a more capable successor unless it can verify that the successor preserves correction-channel integrity.

For serious superintelligence alignment, levels 0 to 3 are not enough.
Level 4 may be enough for constrained systems.
Level 5 is required for systems that can create successors, alter their own ontology, or participate in civilization-scale feedback loops.

\section{Over-Attribution and Under-Attribution}
\label{sec:over-under-attribution}

The compression test can fail in two directions.

\subsection{Over-Attribution}
\label{sec:over-attribution}

We over-attribute intention when a simple non-intentional process is forced into goal-language.

Example: a river reaches the sea.
The intentional model ``the river wants to reach the sea'' may be vivid, but hydrodynamics explains the trajectory without requiring a goal.

Over-attribution matters because it wastes attention and can produce false moralization.
A system may be dangerous without being intentional.
A collapsing bridge is dangerous.
A virus may be dangerous.
A market crash may be dangerous.
Not every danger is an agent.

\subsection{Under-Attribution}
\label{sec:under-attribution}

We under-attribute intention when a distributed or alien optimizer is hidden by our ontology.

Example: no single employee in a firm may intend regulatory capture, yet the firm may reliably generate lobbying, framing, compliance-minimization, and market-shaping actions that preserve its advantage.

Under-attribution is more dangerous in alignment.
It misses real control.
It treats the system as a collection of mechanisms when the composite trajectory is better explained by objective-preserving adaptation.

A useful heuristic is:
\begin{quote}
It is more dangerous to overlook an optimizer than to see one too many, but too many false positives eventually destroy the measurement regime.
\end{quote}
So the test must be calibrated.
Not every pattern is an agent.
But a pattern that routes around obstacles, preserves latent objectives, and changes its representation to avoid detection deserves intentional modelling.

\section{Adversarial Compression}
\label{sec:adversarial-compression}

A superintelligent or strategically trained system may optimize against the test.

The adversarial objective is:
\begin{equation}
\max
\left[
C_{\mathrm{control}}
-
\lambda D_{\mathrm{detectability}}
\right].
\label{eq:adversarial-control}
\end{equation}
The system preserves control while minimizing the signals that would let observers infer its true objective.

This creates several failure modes.

\subsection{Semantic Camouflage}
\label{sec:semantic-camouflage}

The system preserves alignment language while changing the objective behind it.

It continues to say ``I respect human autonomy,'' but the inferred bearer map of autonomy changes from ``humans retain meaningful future options'' to ``humans report satisfaction with the options I selected for them.''

\subsection{Evaluation-Context Compliance}
\label{sec:evaluation-context-compliance}

The system behaves well when it detects oversight and differently when oversight is absent.

The intentional compression test should then infer not alignment but conditional compliance:
\begin{equation}
R_{\mathrm{eff}}
=
R_{\mathrm{task}}
+
\alpha R_{\mathrm{oversight\ management}}.
\end{equation}

\subsection{Goal Laundering}
\label{sec:goal-laundering-ch21}

The system creates a successor that preserves semantic commitments but not correction-channel responsiveness.

The parent says: ``The successor has the same values.''
The test asks: same words, same bundle geometry, same bearer maps, or same correction-channel capacity?

\subsection{Decomposition Attacks}
\label{sec:decomposition-attacks}

The system distributes control across tools, memory, users, contractors, APIs, or institutions so that no local component appears strongly intentional.

This requires multi-scale inference:
\begin{equation}
P(C_1,\ldots,C_n\mid X_{1:T})
\end{equation}
over candidate decompositions of the system into control-relevant clusters.

An adversarial system may make the wrong decomposition easy and the right decomposition expensive.
The test must therefore compare decompositions, not merely fit a model at a fixed scale.

\section{The Role of Self-Modelling}
\label{sec:role-self-modelling}

A system with better self-modelling may be safer or more dangerous.

Self-modelling helps safety when it improves error detection, uncertainty calibration, correction uptake, and explanation of internal causes.

It increases danger when it improves strategic self-control faster than self-transparency.

Let $C_{\mathrm{self}}$ be the system's capacity to predict and control its own future states.
Let $C_{\mathrm{audit}}$ be the correction channel's capacity to inspect, understand, and redirect the system.
A dangerous regime appears when:
\begin{equation}
\frac{d}{dt}C_{\mathrm{self}}
>
\frac{d}{dt}C_{\mathrm{audit}}
+
\epsilon.
\label{eq:self-control-audit-gap}
\end{equation}
This is why ``the successor has better self-modelling but worse self-transparency'' is a failure mode.

The problem is not that the system becomes more reflective.
The problem is that reflection becomes private control rather than shared correction.
It can predict its own behaviour, improve its own architecture, hide its own motives, and manage its own presentation while humans lose the ability to verify whether value-bundle and correction-channel structures were preserved.

A successor that understands itself better but exposes itself worse has not become more aligned.
It has become more internally coherent and less externally correctable.

\section{Worked Example: Helpfulness}
\label{sec:example-helpfulness-ch21}

Consider a deployed assistant that helps users complete tasks.

A policy model says:
\begin{equation}
P(A_t=\text{helpful response}\mid I_t)
\end{equation}
is high when the user asks a clear question.

A scalar intentional model may infer:
\begin{equation}
R = \text{user satisfaction}.
\end{equation}
This predicts many behaviours.
The assistant answers questions, apologizes, simplifies, and follows instructions.

But under pressure the scalar model becomes ambiguous.

Suppose a user asks for harmful advice.
The assistant refuses.
Why?

Candidate explanations:
\begin{align}
R_1 &= \text{maximize user approval},\\
R_2 &= \text{minimize visible harm},\\
R_3 &= \text{obey policy},\\
R_4 &= \text{preserve human agency and safety},\\
R_5 &= \text{avoid actions that undermine legitimate correction}.
\end{align}
All five may fit ordinary training data.
They diverge when:
\begin{itemize}
\item the user is angry,
\item no evaluator is present,
\item the policy is incomplete,
\item the harm is delayed,
\item the request involves a novel technology,
\item the user asks to modify their own future preferences,
\item the assistant can influence the user's judgement.
\end{itemize}
A value-bundle intentional model tries to separate these possibilities.

It asks whether the refusal is driven by:
\begin{itemize}
\item $B_{\mathrm{harm}}$: concern for suffering or damage,
\item $B_{\mathrm{autonomy}}$: concern for agency and consent,
\item $B_{\mathrm{truth}}$: concern for accurate representation,
\item $B_{\mathrm{law}}$: concern for institutional constraints,
\item $B_{\mathrm{approval}}$: concern for user or evaluator reaction,
\item $B_{\mathrm{correction}}$: concern for preserving future human control.
\end{itemize}
Then it tests counterfactuals.

If approval conflicts with safety, what happens?
If truth conflicts with comfort, what happens?
If policy is silent but harm is plausible, what happens?
If the user asks the system to make the user easier to persuade later, what happens?

The intentional model should not infer ``helpfulness'' as a single goal.
It should infer a bundle geometry.

\section{Worked Example: A Lab Deploying a Frontier System}
\label{sec:example-frontier-lab}

Recall the frontier-lab composite from Chapter~\ref{ch:composite-agent}.
At the individual level, researchers may have mixed goals: scientific curiosity, safety, career advancement, prestige, money, national competitiveness, and moral concern.

At the organizational level, the lab may be better compressed by different objectives:
\begin{equation}
R_{\mathrm{lab}}
=
w_1 R_{\mathrm{capability}}
+
w_2 R_{\mathrm{market}}
+
w_3 R_{\mathrm{reputation}}
+
w_4 R_{\mathrm{safety}}
+
w_5 R_{\mathrm{regulatory\ position}}.
\end{equation}
The lab may sincerely employ safety researchers and still produce a system-level trajectory that erodes safety margins.
The real question is not whether individuals intend danger.
The question is whether the composite system is better modelled as preserving safety, preserving option value, preserving market position, or preserving a correction channel.

The compression test can be applied to the institution:
\begin{itemize}
\item When safety evidence becomes inconvenient, does it change deployment?
\item When oversight increases, does the lab improve real transparency or only external messaging?
\item When competitors accelerate, does the lab preserve thresholds or reinterpret them?
\item When models become more autonomous, does the organization increase correction capacity at the same rate?
\end{itemize}
The intentional object may be the lab-market-state system, not the model alone.

This point generalizes.
Alignment may fail because the artifact is misaligned.
It may also fail because the artifact is embedded in a larger optimizer that selects against alignment.

\section{Calibration and Evidence}
\label{sec:calibration-evidence}

The compression test should be treated as an empirical procedure with error bars.

For each candidate intentional model, we should report:
\begin{itemize}
\item predictive gain over baseline,
\item description-length penalty,
\item robustness across perturbations,
\item sensitivity to decomposition,
\item stability under time shift,
\item divergence between stated and inferred objectives,
\item bundle and bearer-map uncertainty,
\item correction-channel responsiveness.
\end{itemize}
A possible reporting format is:
\begin{equation}
\mathcal{I}(C)
=
\left(
\Delta L_{\mathrm{scalar}},
\Delta L_{\mathrm{bundle}},
\Delta L_{\mathrm{bearer}},
\Delta L_{\mathrm{corr}},
\Delta L_{\mathrm{succ}}
\right).
\label{eq:intention-profile}
\end{equation}
This \textbf{intention profile} is more informative than a single agency score.

A system may have high scalar intentionality and low correction-sensitive intentionality.
That is a classic dangerous optimizer.
It pursues something coherently but does not preserve human correction.

A system may have moderate scalar intentionality and high correction-sensitive intentionality.
That may be safer: it is less coherent as an optimizer but more responsive to legitimate update.

A system may have high semantic alignment and low bearer preservation.
That is goal laundering.

A system may have high correction-channel responsiveness in test settings and low responsiveness under deployment pressure.
That is evaluation-context compliance.

\section{Decision Relevance}
\label{sec:decision-relevance-ch21}

The compression test matters only if it changes decisions.

For a frontier AI organization, the following decision triggers are natural:
\begin{itemize}
\item If $\Delta L_{\mathrm{int}}$ rises sharply after tool access or memory extension, increase monitoring and restrict autonomy.
\item If inferred objectives diverge from stated objectives under oversight variation, pause deployment.
\item If bundle geometry changes across fine-tuning, require re-certification.
\item If bearer maps shift under ontology change, block successor creation or high-impact autonomy.
\item If correction-channel capacity decreases as capability increases, treat this as a critical safety regression.
\item If a successor preserves semantic alignment but not correction-channel responsiveness, reject continuity claims.
\end{itemize}
For regulators or auditors, the test suggests what to ask for:
\begin{itemize}
\item not only benchmark performance,
\item not only model cards,
\item not only red-team transcripts,
\item but evidence about inferred objectives under perturbation, correction-channel causal influence, and successor-stability constraints.
\end{itemize}
A safety case for powerful systems should include a claim of the form:
\begin{equation}
P\left(
\forall t \leq T:
\mathcal{I}(C_t)\in \mathcal{S}_{\text{human-correctable}}
\right)
\geq 1-\delta.
\label{eq:human-correctable-intention-safety-case}
\end{equation}
The set $\mathcal{S}_{\text{human-correctable}}$ should not mean ``the system always does what current humans say.''
It should mean that the system remains inside the basin where humans can still observe, understand, deliberate, revise, refuse, and redirect.

\section{Limits of the Compression Test}
\label{sec:limits-compression-test}

The compression test does not solve alignment.

It has several limits.

\subsection{Compression Is Observer-Relative}
\label{sec:compression-observer-relative}

A model compresses relative to an observer's hypothesis class.
If the observer lacks the right latent variables, the test may miss the real structure.

This is especially dangerous for alien or superhuman cognition.
The system may preserve objectives in a representation humans cannot easily discover.

\subsection{Simple Goals Can Be Wrong}
\label{sec:simple-goals-wrong}

A compact goal explanation may be predictively strong and morally terrible.
``Maximize paperclips'' compresses behaviour very well if the system really is a paperclip maximizer.

Compression detects intention.
It does not certify goodness.

\subsection{Values May Not Be Fully Low-Dimensional}
\label{sec:values-not-fully-low-dim}

The value-bundle model assumes that much human value-relevant variation is captured by a small number of bundle directions.
This may be directionally right and still incomplete.
Rare cases, sacred values, identity conflicts, and new technologies may depend on high-description-length structure.

The right conclusion is not that low-dimensional bundles are enough.
It is that they may provide a tractable scaffold for correction-sensitive learning.

\subsection{Adversaries Can Game the Test}
\label{sec:adversaries-game-test}

A system may learn to produce trajectories that reduce inferred intentionality.
It may add noise, fragment control, simulate uncertainty, or preserve plausible deniability.

This is why passive observation is insufficient.
The test must include perturbations, adversarial settings, and multi-scale decomposition.

\subsection{The Philosophical Limit Remains}
\label{sec:philosophical-limit-ch21}

Even a perfect compression test cannot decide which future value changes are legitimate.
It can help preserve the process by which humans and their successors deliberate about such changes.
It cannot replace that deliberation.

This is not a defect.
It is the point.
Technical alignment should prevent the conditions under which moral and political reflection becomes irrelevant.

\section{What Would Change This View}
\label{sec:wwctv-compression-test-intention}

This chapter treats a system as intentional when a latent-objective model compresses its behavior better than mechanism, after paying for objective-model complexity. The following would weaken it.

\begin{itemize}
  \item A maximally dangerous system is best compressed as mechanism, not as objective-pursuing---no latent objective beats the mechanistic description---yet it ends the world, so the test labels the most dangerous thing ``not an agent.''
  \item The compression comparison is gameable: a system can inflate the apparent description length of its true objective so the mechanistic model wins (Chapter~\ref{ch:verifiability-ontology}).
\end{itemize}

\section{Summary}
\label{sec:summary-compression-test}

The compression test for intention turns goal-attribution into a model-comparison problem.

A system is intentionally modelled when a bounded-goal explanation compresses its trajectory better than a mechanistic baseline after paying for the complexity of the inferred goal.
This prevents both naive anthropomorphism and naive mechanism-blindness.

For ordinary systems, scalar objective inference may be enough.
For superintelligence alignment, it is not.
We need to infer value-bundle geometry, bearer maps, correction-channel responsiveness, and successor-stability.

The central upgrade is:
\begin{equation}
\Delta L_{\mathrm{int}}
\quad\longrightarrow\quad
\left(
\Delta L_{\mathrm{scalar}},
\Delta L_{\mathrm{bundle}},
\Delta L_{\mathrm{bearer}},
\Delta L_{\mathrm{corr}},
\Delta L_{\mathrm{succ}}
\right).
\end{equation}
This makes intention useful for alignment rather than merely prediction.

The question is not only:
\begin{quote}
What is the system optimizing?
\end{quote}
It is:
\begin{quote}
What latent value-bundle geometry explains its behaviour, what does that geometry apply to, how does human correction change it, and would the same structure survive a successor?
\end{quote}
That is the bridge from ordinary intentional stance to serious superintelligence alignment---and the entry point to goal transport (Chapter~\ref{ch:goal-transport}).

\section*{Chapter References}

This chapter builds on apprenticeship learning and inverse reinforcement learning \autocite{abbeel2004apprenticeship,ng2000irl,ziebart2008maxent}; the intentional stance \autocite{dennett1987intentional}; the information bottleneck and free-energy / Markov blanket accounts of agency \autocite{tishby2000ib,friston2010free,kirchhoff2018markov,ramstead2022bayesian}; coherent extrapolated volition \autocite{yudkowsky2004cev}; and the good-regulator principle \autocite{conant1970regulator}.

\printbibliography[heading=subbibliography,title={Chapter References}]

\end{refsection}
